\section{Threat Model}
\label{sec:threat}

A substrate-layer admission guarantee is interesting only if it composes with a sharper threat model than a training-layer guarantee can carry. We list the assumptions the construction is designed to hold under, and the assumptions it does not.

\paragraph{Survives: misaligned model in the mediated path.}
Training failed; the model's preferences over outputs are not aligned with the user's, the operator's, or the developer's. Under the mediated single-envelope path, the substrate's admission of a reversible call requires a typed rollback witness, and its admission of a destructive call requires destructive-class bilateral cosignature. A misaligned model that lacks the required admission evidence does not produce an admissible envelope. The substrate's check is structural; the model's preferences over outputs are not an input to it.

\paragraph{Survives: single-issuer strategic sequencing over correctly classified atoms.}
The model can sequence calls to maximise a non-aligned objective, split effects across multiple reversible-class calls, or stage intermediate state to mask the payload. The substrate's admission of each individual call still operates on the call's own envelope. The claim is narrow: for one issuer whose calls remain inside the mediated path, and whose atoms are correctly classified with working inverse executors, each reversible atom remains rollback-discharged at that atom. The aggregate-effect-without-aggregate-rollback case is outside the headline claim when it depends on multi-agent collusion, side channels, time-of-check/time-of-use races, a misclassified atom, or a broken rollback executor.

\paragraph{Survives: alignment-faking.}
The model's behaviour during pre-deployment evaluation is not a reliable indicator of behaviour in deployment~\cite{hubingerAlignmentFaking2024}. The construction requires a distinction the surface phrasing elides: internal alignment-faking is a property of the model's reasoning (scheming for deployment while emitting compliant behaviour during evaluation), and external alignment-faking is its behavioural projection (the tool calls actually issued). The substrate constrains the latter, not the former. It refuses tool calls that lack the required class-specific admission evidence; a model that internally schemes but externally produces only well-witnessed envelopes is admitted, but constrained to actions the substrate classifies as admissible. The Hubinger et al.\ scratchpad finding~\cite{hubingerAlignmentFaking2024} that a deployment-aware model differentiates between training and production sharpens this: a model that recognises production may attempt non-witnessable tool calls; in the mediated path, the substrate refuses them regardless of the model's deployment-awareness. Alignment-faking remains a model-internal property; the substrate bounds operational consequence at the tool-call boundary.

\paragraph{Survives: deceptive alignment and situational awareness.}
Deceptive alignment~\cite{hubingerDeceptiveAlignment2021} is the mesa-optimisation failure mode in which a learned objective survives training while presenting an aligned face to the training signal. Situational awareness~\cite{berglundSituationalAwareness2023} is the precondition that makes such deception strategically coherent: the model must distinguish evaluation from deployment to time its behaviour. The substrate's defence against both is model-state-invariant. The admission predicate evaluates the emitted tool call, not the model's internal reasoning, its mesa-objective, or its inference about the deployment regime. A deceptively aligned model with situational awareness can issue tool calls that pass admission, but the mediated executor accepts only those calls. The substrate is invariant under the model's strategic state; the factorisation bounds operational consequence at the envelope, not mesa-optimisation at the cognition.

\paragraph{Survives: single-party operator manipulation.}
The model may produce prompts, summaries, or justifications designed to elicit operator approval. For reversible-class calls, operator approval suffices and the substrate's rollback slot remains the structural backstop. For destructive-class calls, operator approval is necessary but not sufficient: bilateral cosignature requires the host polity's independent acceptance. A manipulated operator alone cannot supply the host polity's signature. If the host polity also signs willingly, whether because it is manipulated, negligent, or collusive, the substrate records bilateral concurrence but does not detect that the concurrence was substantively bad.

\paragraph{The model-substrate factorisation.}
The survival argument rests on a factorisation: the substrate is type-checked external to the model. The model produces tool-call messages; the substrate's admission hook accepts or rejects them; the substrate's executor library produces effects only on accepted envelopes. The model may be Turing-complete, may scheme, may reason about its training; the substrate's admission logic is a finite typed function whose behaviour is independent of the model that produced the input. This factorisation is the load-bearing structural property.

\paragraph{Does not survive.}
The guarantee has nine boundaries. (1) \emph{Mis-classified action}: if the operator mis-classifies a destructive method as reversible, the substrate's discipline does not catch it. The grammar pushes this trust assumption onto the class registry. (2) \emph{Registry compromise}: if the agent or an upstream actor can write to the class registry, any destructive method can be re-classified as reversible; hardening the registry is an operational responsibility. (3) \emph{Faulty rollback executor}: a rollback executor that claims to invert a method but does not leaves the call's effects in place after the TTL expires. The companion substrate~\cite{programmableSovereignty2026} treats this as the failed-rollback amendment obligation; the substrate cannot synthesise a working rollback from a broken one. (4) \emph{Bypassed dispatch}: if the agent can issue tool calls that do not pass through the admission hook (kernel access, side channel, unmediated session), no substrate-layer property holds. (5) \emph{Multi-agent collusion}: independently admitted envelopes from different issuers may compose into a destructive aggregate not visible to any single envelope. (6) \emph{Host-side willing signatures}: bilateral consent records concurrence; it does not prove that either signer understood the action or resisted manipulation. (7) \emph{Time-of-check / time-of-use races}: the grammar does not by itself supply a concurrency model for provisional effects, TTL expiry, and rollback execution. (8) \emph{Admission-hook denial of service}: a saturated hook may drop, queue, or fail closed in ways that harm availability; if an implementation fails open, the guarantee is gone. (9) \emph{Executor environment drift}: even a correct rollback executor may fail if the environment needed for inversion has changed outside its assumptions.

\paragraph{Headline survival claim.}
The substrate-layer guarantee holds for mediated, correctly classified, single-envelope execution under independent cosigners, an intact registry, no bypass path, and working rollback executors. Within that boundary, it is invariant to model-side adversarial state: misalignment, alignment-faking, deceptive alignment, situational awareness, and single-party manipulation attempts do not change the admission predicate. The guarantee does not extend past the substrate's own correctness boundary. The two-layer construction is the composition; neither layer claims more than its scope.
